\section{مجموعه‌های ۲ تا ۴: مقایسه جفت‌ها}
\label{sec:sets234}
در این بخش فقط داده نهایی (بعد از همه برش‌ها، با جت‌های کامل و جزئی با هم) استفاده می‌شود. هر مجموعه دو نمونه را مقایسه می‌کند و زیرنمودار نسبت دارد. ۱۶ متغیر در چهار شکل چهارتایی آمده‌اند. ترتیب متغیرها در هر مجموعه یکی است.

\subsection{مجموعه ۲: \ttZ در برابر \lr{Z+jet} (دو نمونه سیگنال)}
\label{sec:set2}
\quadfig{set2_ttZ_vs_Zj_pt.pdf}{set2_ttZ_vs_Zj_eta.pdf}{set2_ttZ_vs_Zj_mass.pdf}{set2_ttZ_vs_Zj_msd.pdf}{مجموعه ۲: \pt، $\eta$، جرم و جرم \lr{soft-drop}.}{fig:s2a}
\quadfig{set2_ttZ_vs_Zj_girth.pdf}{set2_ttZ_vs_Zj_nconst.pdf}{set2_ttZ_vs_Zj_tau1.pdf}{set2_ttZ_vs_Zj_tau2.pdf}{مجموعه ۲: \lr{girth}، \lr{nconst}، $\tau_1$ و $\tau_2$.}{fig:s2b}
\quadfig{set2_ttZ_vs_Zj_tau3.pdf}{set2_ttZ_vs_Zj_tau21.pdf}{set2_ttZ_vs_Zj_tau32.pdf}{set2_ttZ_vs_Zj_c2.pdf}{مجموعه ۲: $\tau_3$، $\tau_{21}$، $\tau_{32}$ و $C_2$.}{fig:s2c}
\quadfig{set2_ttZ_vs_Zj_d2.pdf}{set2_ttZ_vs_Zj_n2.pdf}{set2_ttZ_vs_Zj_zg.pdf}{set2_ttZ_vs_Zj_rg.pdf}{مجموعه ۲: $D_2$، $N_2$، $z_g$ و $R_g$.}{fig:s2d}

دو نمونه سیگنال شکل بسیار مشابهی دارند:
\begin{itemize}
\item \textbf{جرم و جرم \lr{soft-drop}:} قله هر دو نزدیک ۹۰~\lr{GeV} است، یعنی در هر دو نمونه جت کامل \lr{Z} را بازسازی کرده‌ایم. در \ttZ قله کمی پایین‌تر و دم جرم بالا بلندتر است (نسبت تا ۲ و ۳).
\item \textbf{$\eta$:} \ttZ مرکزی‌تر است. نسبت \ttZ به \lr{Z+jet} در $\eta=0$ حدود ۱.۴ و در لبه $|\eta|\approx 2$ حدود ۰.۳ است.
\item \textbf{\pt:} یکسان تا ۶۰۰ \lr{GeV}؛ بعد از آن \ttZ کمی سخت‌تر است.
\item \textbf{\lr{girth}، $\tau_N$، $C_2$، \lr{nconst}:} در \ttZ توزیع‌ها کمی به مقدار بزرگ‌تر جابه‌جا شده و دم بلندتر دارند. علت‌های محتمل در بخش~\ref{sec:zorigin} بررسی شده است.
\item \textbf{$\tau_{21}$، $N_2$، $R_g$:} تقریباً یکسان‌اند (نسبت نزدیک ۱).
\item \textbf{$\tau_{32}$:} در مقدار کم $\tau_{32}$ (کمتر از ۰.۵) \ttZ دو برابر است؛ بقیه یکسان.
\item \textbf{$D_2$:} \ttZ در $D_2$ کوچک کمی بیشتر و در $D_2$ بزرگ کمتر است (نسبت ۰.۵).
\end{itemize}
نتیجه: دو نمونه سیگنال از نظر ساختار دو‌هسته‌ای هم‌خوانی دارند و تفاوت‌ها را در بخش~\ref{sec:zorigin} با تفاوت $\eta$، قطبش \lr{Z} و آلودگی ناشی از واپاشی تاپ توضیح می‌دهیم.

\subsection{مجموعه ۳: \ttbb در برابر \lr{bbbb} (دو نمونه پس‌زمینه)}
\label{sec:set3}
\quadfig{set3_ttbb_vs_bbbb_pt.pdf}{set3_ttbb_vs_bbbb_eta.pdf}{set3_ttbb_vs_bbbb_mass.pdf}{set3_ttbb_vs_bbbb_msd.pdf}{مجموعه ۳: \pt، $\eta$، جرم و جرم \lr{soft-drop}.}{fig:s3a}
\quadfig{set3_ttbb_vs_bbbb_girth.pdf}{set3_ttbb_vs_bbbb_nconst.pdf}{set3_ttbb_vs_bbbb_tau1.pdf}{set3_ttbb_vs_bbbb_tau2.pdf}{مجموعه ۳: \lr{girth}، \lr{nconst}، $\tau_1$ و $\tau_2$.}{fig:s3b}
\quadfig{set3_ttbb_vs_bbbb_tau3.pdf}{set3_ttbb_vs_bbbb_tau21.pdf}{set3_ttbb_vs_bbbb_tau32.pdf}{set3_ttbb_vs_bbbb_c2.pdf}{مجموعه ۳: $\tau_3$، $\tau_{21}$، $\tau_{32}$ و $C_2$.}{fig:s3c}
\quadfig{set3_ttbb_vs_bbbb_d2.pdf}{set3_ttbb_vs_bbbb_n2.pdf}{set3_ttbb_vs_bbbb_zg.pdf}{set3_ttbb_vs_bbbb_rg.pdf}{مجموعه ۳: $D_2$، $N_2$، $z_g$ و $R_g$.}{fig:s3d}

دو نمونه پس‌زمینه تقریباً یکسان‌اند. نسبت \ttbb به \lr{bbbb} برای همه متغیرهای زیرساختار نزدیک ۱ است. تنها تفاوت‌ها این‌هاست:
\begin{itemize}
\item \pt: \ttbb دم سخت‌تری دارد (نسبت ۲ تا ۳ در \pt بالای ۸۰۰~\lr{GeV})؛ دلیلش برش پارتونی $\pt(\mathrm{b\bar{b}})>400$ است که فقط روی \ttbb است.
\item جرم و \lr{nconst}: در دم بالا \ttbb تا ۱.۵ برابر است.
\item $\eta$: یکسان.
\end{itemize}
نتیجه: دو پس‌زمینه از نظر زیرساختار قابل ادغام‌اند.

\subsection{مجموعه ۴: سیگنال در برابر پس‌زمینه}
\label{sec:set4}
در این مجموعه دو جفت (\lr{Z+jet} در برابر \lr{bbbb}) و (\ttZ در برابر \ttbb) با هم رسم شده‌اند. تفاوت دو جفت بیشتر درباره انتقال‌پذیری اهمیت دارد.
\quadfig{set4_signal_vs_background_pt.pdf}{set4_signal_vs_background_eta.pdf}{set4_signal_vs_background_mass.pdf}{set4_signal_vs_background_msd.pdf}{مجموعه ۴: \pt، $\eta$، جرم و جرم \lr{soft-drop}.}{fig:s4a}
\quadfig{set4_signal_vs_background_girth.pdf}{set4_signal_vs_background_nconst.pdf}{set4_signal_vs_background_tau1.pdf}{set4_signal_vs_background_tau2.pdf}{مجموعه ۴: \lr{girth}، \lr{nconst}، $\tau_1$ و $\tau_2$.}{fig:s4b}
\quadfig{set4_signal_vs_background_tau3.pdf}{set4_signal_vs_background_tau21.pdf}{set4_signal_vs_background_tau32.pdf}{set4_signal_vs_background_c2.pdf}{مجموعه ۴: $\tau_3$، $\tau_{21}$، $\tau_{32}$ و $C_2$.}{fig:s4c}
\quadfig{set4_signal_vs_background_d2.pdf}{set4_signal_vs_background_n2.pdf}{set4_signal_vs_background_zg.pdf}{set4_signal_vs_background_rg.pdf}{مجموعه ۴: $D_2$، $N_2$، $z_g$ و $R_g$.}{fig:s4d}

\subsubsection*{متغیرهای جداکننده}
\begin{itemize}
\item \textbf{جرم و $m_{sd}$:} سیگنال قله تیز نزدیک ۹۰~\lr{GeV}؛ پس‌زمینه پهن و بدون قله. نسبت سیگنال به پس‌زمینه در ناحیه قله ۱.۸ تا ۲.۳ است و زیر ۱ در بقیه جاها.
\item \textbf{\lr{girth}، $\tau_1$:} سیگنال در ۰.۱۵ تا ۰.۲۵ قله دارد؛ پس‌زمینه در مقادیر کمتر و دم بلندتر. نسبت تا ۲.۵ و زیر ۱ در مقادیر کوچک و بزرگ.
\item \textbf{$\tau_{21}$ و $C_2$ و $D_2$ و $N_2$:} سیگنال دو‌هسته‌ای‌تر است ($\tau_{21}$ و $C_2$ و $D_2$ کوچک‌تر، $N_2$ کوچک‌تر).
\item \textbf{$\tau_2$ و $\tau_3$:} سیگنال در مقدار کم متمرکزتر است؛ اثر در $\tau_3$ ضعیف‌تر است.
\item \textbf{\lr{nconst}:} سیگنال مؤلفه کمتر دارد (قله حدود ۴۰ در برابر ۴۷). نسبت تا ۳ در تعداد کم و ۰.۲ در بیش از ۸۰ مؤلفه.
\item \textbf{$R_g$:} سیگنال قله در ۰.۳ تا ۰.۴ دارد (دو زیرجت با فاصله مشخص)، پس‌زمینه در مقدار کم (زیرجت‌ها نزدیک هم) و دم تخت. نسبت تا ۲.۸.
\end{itemize}

\subsubsection*{متغیرهای کم‌اثر یا وابسته به جفت}
\begin{itemize}
\item \textbf{$\tau_{32}$:} تقریباً هیچ تفاوتی بین سیگنال و پس‌زمینه نیست (نسبت $\approx$۱). فقط در $\tau_{32}$ کم، \ttZ/\ttbb بالای ۱ و \lr{Z+jet}/\lr{bbbb} زیر ۱ است.
\item \textbf{$z_g$:} نسبت بین ۰.۶ و ۱.۵ و شیب ملایم؛ جداکنندگی ضعیف و کمی متفاوت در دو جفت.
\item \textbf{\pt:} شکل مشابه؛ نسبت تقریباً ۱ تا حدود ۷۰۰~\lr{GeV} و بعد پراکنده.
\item \textbf{$\eta$:} دو جفت رفتار \emph{برعکس} دارند. \lr{Z+jet} از \lr{bbbb} پهن‌تر است (نسبت ۱.۵ در $|\eta|\approx 2$ و ۰.۸ در مرکز)، ولی \ttZ از \ttbb مرکزی‌تر است (نسبت ۰.۴ در لبه و ۱.۲ در مرکز).
\end{itemize}

\subsubsection*{مشاهده مهم}
برای جداکننده‌های اصلی (جرم، $m_{sd}$، \lr{girth}، $\tau_N$، $\tau_{21}$، $C_2$، $D_2$، $N_2$، $R_g$، \lr{nconst}) نسبت سیگنال به پس‌زمینه برای هر دو جفت تقریباً روی هم می‌افتد؛ یعنی جداکنندگی با تغییر جفت تغییر نمی‌کند. فقط $\eta$ (و در درجه کمتر \pt و $z_g$ و $\tau_{32}$ در ناحیه‌های خاص) بین دو جفت متفاوت است.


\section{منشأ تفاوت \lr{Z} در \lr{Z+jet} و \ttZ}
\label{sec:zorigin}
در مجموعه ۲ دیدیم که جت‌های \Zbb در \lr{Z+jet} و \ttZ در $\eta$، جرم، \lr{girth} و چند متغیر دیگر با هم فرق دارند. اگر بوزون \lr{Z} در دو فرایند دقیقاً یکی بود، انتظار می‌رفت جتش هم یکی باشد، پس باید دید این تفاوت از کدام مرحله شروع می‌شود. برای همین خروجی هر مرحله شبیه‌سازی را نگه داشته‌ایم و \lr{Z} را در خروجی \lr{MadGraph} (پیش از واپاشی) و \lr{MadSpin} (پس از واپاشی) برای ۲۰۰۰۰ رویداد هر نمونه مقایسه می‌کنیم.

نخست خود تولید \lr{Z}. توزیع \pt در دو نمونه تقریباً یکسان است (میانگین ۴۹۹ در \lr{Z+jet} و ۵۰۹~\lr{GeV} در \ttZ، چون در هر دو $\pt(\mathrm{Z})>400$~\lr{GeV} گذاشته‌ایم)؛ دم \ttZ کمی سخت‌تر است و سهم $\pt>700$~\lr{GeV} برابر ۶.۷\% در برابر ۵.۶\% است. ولی $|\eta|$ و $|y|$ در \ttZ کوچک‌ترند: میانگین $|\eta|$ برابر ۰.۹۶ در \lr{Z+jet} و ۰.۶۹ در \ttZ است (شکل~\ref{fig:lhe-prod}). دلیل این تفاوت آستانه انرژی فرایند سخت است. اگر دو پارتون ورودی کسرهای $x_1$ و $x_2$ از تکانه پروتون را داشته باشند، انرژی فرایند سخت و تندی سیستم برابرند با
\[
\hat s = x_1 x_2\, S,\qquad y=\frac12\ln\frac{x_1}{x_2},
\]
که $S=(13~\mathrm{TeV})^2$ است. از اینجا $x_{1,2}=\sqrt{\hat s/S}\,e^{\pm y}$ و چون باید $x\le1$ باشد،
\[
|y|\le \ln\frac{\sqrt{S}}{\sqrt{\hat s}}.
\]
یعنی هرچه $\hat s$ بزرگ‌تر باشد، بازه مجاز تندی کوچک‌تر می‌شود. در \ttZ دستگاه $t\bar t$ هم، که جرم دارد، باید تولید و پس‌زده شود و آستانه $\hat s$ بالاتر می‌رود. داده همین را نشان می‌دهد: میانه $\sqrt{\hat s}$، یعنی جرم ناوردای همه ذرات خروجی فرایند سخت، در \ttZ برابر ۱۴۹۲ و در \lr{Z+jet} برابر ۱۰۴۸~\lr{GeV} است و میانگین $|y|$ کل سیستم ۰.۵۵ در برابر ۰.۸۰؛ در پس‌زمینه‌ها هم همین جهت دیده می‌شود (میانه $\sqrt{\hat s}$ برابر ۱۵۱۲ در \ttbb و ۱۱۰۶ در \lr{bbbb}، و میانگین $|y|$ برابر ۰.۵۲ و ۰.۶۱). چون $\eta$ جت از تندی سیستم می‌آید، سیستم باریک‌تر یعنی جت مرکزی‌تر.

\quadfig{lhe_Z_pt.pdf}{lhe_Z_abs_eta.pdf}{lhe_Z_abs_y.pdf}{lhe_Z_mass.pdf}{توزیع \pt، $|\eta|$، $|y|$ و جرم \lr{Z} در خروجی \lr{MadGraph} (خط‌چین) و \lr{MadSpin} (خط پر) برای \lr{Z+jet} و \ttZ. منحنی‌های قبل و بعد از \lr{MadSpin} روی هم افتاده‌اند و جرم \lr{Z} روی ۹۱.۱۸۸~\lr{GeV} ثابت است.}{fig:lhe-prod}

واپاشی با \lr{MadSpin} تولید \lr{Z} را عوض نمی‌کند: میانگین \pt و $|\eta|$ قبل و بعد از آن در هر دو نمونه تا سه رقم اعشار یکی است و جرم \lr{Z} در هر دو مرحله روی ۹۱.۱۸۸~\lr{GeV} می‌ماند، چون \lr{MadSpin} در حالت \lr{onshell} اجرا می‌شود. پس پهنای \lr{Z} در این تولید اثری ندارد و تفاوت $\eta$ و دم \pt از خود ماتریس‌المان در \lr{MadGraph} می‌آید.

ولی خود واپاشی $\mathrm{Z}\to\mathrm{b}\bar{\mathrm{b}}$ در دو نمونه یکی نیست (شکل~\ref{fig:lhe-decay}). زاویه $\theta^*$ را زاویه کوارک \lr{b} در دستگاه سکون \lr{Z} نسبت به جهت حرکت \lr{Z} در آزمایشگاه می‌گیریم. برای \lr{Z} عرضی توزیع $\cos\theta^*$ شبیه $1+\cos^2\theta^*$ و برای \lr{Z} طولی شبیه $1-\cos^2\theta^*$ است. اگر فقط این دو حالت را در نظر بگیریم، $\langle\cos^2\theta^*\rangle=0.4-0.2f_L$ و کسر طولی از $f_L=2-5\langle\cos^2\theta^*\rangle$ به دست می‌آید. نتیجه $f_L=0.07\pm0.01$ در \lr{Z+jet} و $0.57\pm0.01$ در \ttZ است (خطای آماری)، و در هر سه بازه \pt (۴۰۰ تا ۴۵۰، ۴۵۰ تا ۵۵۰ و بالاتر) همین الگو دیده می‌شود. یعنی \lr{Z} در \lr{Z+jet} تقریباً عرضی است و در \ttZ بیش از نیمی از آن طولی است. در \lr{Z+jet} میانگین $\cos\theta^*$ هم حدود ۰.۱۳ است، یعنی توزیع نسبت به b و $\bar{\mathrm b}$ متقارن نیست؛ این جمله فرد در $\langle\cos^2\theta^*\rangle$ اثری ندارد و آن را بررسی نکرده‌ایم. \textcolor{red}{علت فیزیکی تفاوت قطبش در دو فرایند را هنوز بررسی نکرده‌ایم.}

\quadfig{lhe_dR_bb.pdf}{lhe_pt_b_sub.pdf}{lhe_cos_theta_star.pdf}{lhe_sqrt_z1mz.pdf}{واپاشی \lr{Z} پس از \lr{MadSpin}: $\Delta R(\mathrm{b},\bar{\mathrm b})$، \pt کوارک \lr{b} دوم، $\cos\theta^*$ (با کسر طولی $f_L$) و $\sqrt{z(1-z)}$ که $z$ سهم \lr{b} نرم‌تر از مجموع \pt دو \lr{b} است.}{fig:lhe-decay}

قطبش تقسیم انرژی بین دو \lr{b} را عوض می‌کند: \lr{Z} طولی تقسیم متقارن‌تری می‌دهد. میانگین \pt کوارک \lr{b} دوم برابر ۱۱۴ در \lr{Z+jet} و ۱۴۰~\lr{GeV} در \ttZ است و میانگین $\sqrt{z(1-z)}$ برابر ۰.۳۷۶ در برابر ۰.۴۱۲ می‌شود. اگر \lr{Z} به دو زیرجت با سهم $z$ و $1-z$ از تکانه بشکند، $\Delta R(\mathrm{b},\bar{\mathrm b})\approx m/(\pt\sqrt{z(1-z)})$ و \lr{girth} حدود $2\sqrt{z(1-z)}\,m/\pt$ است، پس تقسیم متقارن‌تر $\Delta R$ را کوچک‌تر و \lr{girth} را بزرگ‌تر می‌کند. در داده هم همین دیده می‌شود: میانگین $\Delta R(\mathrm{b},\bar{\mathrm b})$ برابر ۰.۵۹ در \lr{Z+jet} و ۰.۵۰ در \ttZ است، ولی میانه \lr{girth} جت‌های کامل ۰.۱۴۵ در \lr{Z+jet} و ۰.۱۵۹ در \ttZ است. نسبت \ttZ به \lr{Z+jet} در \lr{girth} حدود ۱.۱۰ است و پیش‌بینی ساده از $\langle\sqrt{z(1-z)}\rangle$ هم ۱.۱۰ (۰.۴۱۲ به ۰.۳۷۶) می‌دهد. این هم‌خوانی تقریبی است، چون فرمول‌ها تابش و ذره‌های نرم را در نظر نمی‌گیرند.

منبع دوم تفاوت آلودگی ناشی از واپاشی تاپ است. در \ttZ، ۸.۳\% از جت‌های کامل دست‌کم یک \lr{b} از تاپ دارند و در \lr{Z+jet} هیچ‌کدام. در سطح ذره‌های \lr{Pythia} (۴۰۰۰ رویداد هر نمونه)، در ۱۴.۷\% رویدادهای \ttZ دست‌کم یک ذره از نسل تاپ داخل مخروط $\Delta R<0.8$ دور \lr{Z} است؛ این نسل را از روی اولین مادر هر ذره تشخیص داده‌ایم و بنابراین تقریبی است. این ذره‌ها دم جرم را بالا می‌برند: میانگین جرم ذره‌های مخروط ۱۰۱.۹ در \ttZ و ۹۰.۳~\lr{GeV} در \lr{Z+jet} است و اگر فقط ذره‌های نسل \lr{Z} را بشماریم ۸۱.۳ و ۷۸.۳ می‌شود. در جت‌های نهایی هم حذف همان ۸.۳\% جت‌ها میانه جرم \ttZ را از ۹۳.۱ به ۹۱.۹ (\lr{Z+jet}: ۹۰.۱)، میانه \lr{nconst} را از ۴۳ به ۴۲ (۴۱) و میانه \lr{girth} را از ۰.۱۵۹ به ۰.۱۵۵ (۰.۱۴۵) می‌رساند، یعنی نسبت \lr{girth} از ۱.۱۰ به ۱.۰۷ کم می‌شود. \textcolor{red}{بعد از حذف این جت‌ها جرم و \lr{nconst} در \ttZ هنوز حدود ۲\% بیشتر است؛ علت این باقی‌مانده را نمی‌دانیم.}

در مجموع، \lr{Z} در \lr{Z+jet} و \ttZ از همان مرحله تولید فرق دارد (در $\eta$ و قطبش) و بخشی از تفاوت جت‌ها فیزیکی است، نه خطای شبیه‌سازی. این یعنی شناساگری که فقط روی یکی از دو نمونه آموزش دیده با \lr{Z}ای روبه‌رو می‌شود که در ساختار دو‌هسته‌ای هم با آن فرق دارد، و این را باید در آزمون انتقال‌پذیری در نظر گرفت.
